% The optional write-up, Section 4.2. Build with:
%     latexmk -pdf write_up.tex && latexmk -c
% which leaves write_up.pdf beside this file, where src/submission/package.py
% looks for it (src.common.config.WRITE_UP_DIR). Sources, figures and the PDF
% stay together so the packager needs no build step of its own.
\documentclass[11pt,a4paper]{article}

\usepackage[T1]{fontenc}
\usepackage[utf8]{inputenc}
\usepackage[margin=2cm]{geometry}
\usepackage{amsmath}
\usepackage{booktabs}
\usepackage{enumitem}
\usepackage{parskip}
\usepackage{microtype}
\usepackage{titlesec}
\usepackage[hidelinks]{hyperref}

\setlist{nosep,leftmargin=1.2em}
% Paths and identifiers set in \texttt do not hyphenate, so give TeX room to
% stretch a line rather than push one into the margin.
\emergencystretch=3em
% Three pages for four sentences of overflow is not a short write-up. Tightening the
% headings buys the last two pages back without shrinking any of the text.
\titlespacing*{\section}{0pt}{1.1ex plus .2ex}{0.6ex}
\newcommand{\app}{\textsc{neb-x}}

\title{\vspace{-1.4cm}\app: Train Condition Monitoring\\[0.2em]
\large Problem Statement 3, all four subsystems}
\author{}
\date{}

\begin{document}
\maketitle
\vspace{-1.6cm}

\section{What we built}

One web app that a depot engineer can use without touching code. Upload the recordings from a
system, get an answer, read why, and download a maintenance handoff package. It covers all four
subsystems, and it is the same app that produced every row of our \texttt{predictions.zip}.

\begin{center}\small
\begin{tabular}{@{}llll@{}}
\toprule
\textbf{Subsystem} & \textbf{Task} & \textbf{Metric} & \textbf{Our validation figure} \\
\midrule
Door  & Segment cycles, classify each   & IoU-weighted F1 & 1.000 accuracy, 110 cycles \\
ACV   & Rank cars by leak likelihood    & Linear rank-decay & faulty car first, 5 of 5 cases \\
Rail  & Normal / Side I / Side II       & Macro F1 & 0.750 $\pm$ 0.114, 0.757 pooled \\
SHM   & Cumulative fatigue damage       & $\max(0, 1-\mathrm{MAPE})$ & MAPE 2.52\%, score 0.9748 \\
\bottomrule
\end{tabular}
\end{center}

Live at \url{https://neb-x-357243228666.asia-southeast1.run.app}; \texttt{./scripts/app.sh} runs
it locally.

\section{The app and its pipeline}

Six steps, in the order a user meets them:

\begin{enumerate}
  \item \textbf{Choose a system.} A board shows the four systems and whether each has a model
        loaded, so an unavailable one says so rather than failing on the first click.
  \item \textbf{Add recordings.} Drag and drop; each system states the file types it reads.
        Filenames are preserved verbatim, because \texttt{file\_id} in the submission is the
        source filename and a reformatted one matches nothing.
  \item \textbf{Assess.} The batch runs in the background, one system at a time, so the user can
        switch to another system while the first finishes.
  \item \textbf{Read the answer.} A verdict panel gives the headline, a severity colour and one
        line of detail. \emph{Answer} is the plain-language view; \emph{Details} steps through the
        model's own workings one panel at a time: per-cycle strips, ranked bars, threshold rows,
        time series.
  \item \textbf{Keep the session.} Every system keeps its uploads, its latest assessment and its
        selected cycle while the user moves between systems.
  \item \textbf{Hand off.} One button builds a ZIP: a plain-text summary, an offline printable
        HTML report with charts, per-system findings as CSV with evidence charts, and a
        \texttt{manifest.json} carrying assessment times, source filenames and checksums. No raw
        recordings, no model binaries, and systems with nothing to report are documented as such.
\end{enumerate}

Every explanation the app draws is computed by the same functions that build the feature vector,
so an explanation cannot drift from the prediction it explains.

\section{How each subsystem is trained}

\textbf{Door.} The controller only samples while a door moves, so cycles are separated by silence:
rows are 20\,ms apart inside a cycle and the recorder goes quiet for 10\,s or more between them.
Splitting on that gap reproduces all 110 labelled cycles exactly, which makes every IoU in the
official metric 1.0 and collapses the score to plain accuracy. Classification then uses one
feature: a cycle's mean motor current divided by the median of its own operation \emph{within its
own recording}. Absolute milliamps do not transfer between doors (info kit \S1.2.2); a cycle's
excess over its own door's typical cycle does. Scaling all currents by 0.85--1.20$\times$ costs a
raw threshold up to 80 of 110 cycles and costs this feature nothing. The estimator is a gradient
boosting classifier, effectively a stump that derives its own split, kept over a hardcoded
threshold so the app has a probability to show. Validation is repeated stratified 5-fold,
10 repeats.

\textbf{ACV.} Nothing is fitted. Each car reports its cabin temperature and its cooling setpoint;
subtract, average over the file, rank descending. A healthy car sits at or below its target, a
leaking car runs above its own target because it cannot keep up. That single signal ranks the true
faulty car first in all five cases sharing the held-out file's schema. Every car is always ranked,
never dropped: the organisers score an omitted car as zero.

\textbf{Rail Corrugation.} One file is one second of an 8-car train crossing a section of track,
64 axle boxes per rail at 10\,kHz. Both rails are crossed by the same train at the same speed in
the same second, so the \emph{contrast between the two sides} cancels the confounders the info kit
names (speed, ballast noise, track elasticity) and what remains is genuinely asymmetric. Energy is
binned by ripple \emph{wavelength} rather than frequency, because a wheel crossing corrugation
rings at $\text{speed}/\text{wavelength}$ and speed varies tenfold across these files; that choice
alone is worth $+0.06$ macro F1. Each side's 32 boxes are aggregated by max and p90, never median,
because only some boxes cross the corrugated stretch and a median is built to discard exactly
those. Estimator: class-balanced gradient boosting, 228 features, repeated stratified 5-fold with
10 repeats. A regularised linear baseline ties it, and five hyperparameter variations fail to beat
it, so the simple model is the one we ship.

\textbf{SHM.} Physics, with two fitted numbers. Rainflow counting reduces each stress history to
closed cycles, and Miner's rule sums them into damage $D = S(m)/C$ with
$S(m)=\sum_i n_i \sigma_i^{m}$. The exponent $m$ and coefficient $C$ are chosen to minimise MAPE
directly, the metric this subsystem is scored on, rather than least squares, which would chase the
large-damage files and leave the small ones wrong. The search is coarse then fine because the
optimum in $m$ is sharp. Validation is leave-one-out over the 64 labelled files, honest here only
because the model has two parameters; with a high-capacity model it would leak.

\section{How it comes together}

Dependencies point one way: \texttt{common} $\leftarrow$ subsystem $\leftarrow$ app. No subsystem
imports another, shared logic moves down into \texttt{common/} and never sideways, and every path
in the project is owned by one file (\texttt{src/common/config.py}) so the trainer and the
predictor cannot disagree about where a checkpoint lives.

The app never branches on which subsystem is selected. It asks for a capability and gets one or
does not: \texttt{services.py} resolves each subsystem's \texttt{predict} (required) and
\texttt{explain} (optional) by \texttt{importlib}, over the same input paths. Explanations come back
as plain dictionaries tagged with a \texttt{kind}, and the renderer draws whatever arrives. A fifth
subsystem would be added without editing a line of the app.

The submission is one command. \texttt{./submit.sh} runs every held-out test input through the
same \texttt{predict} the app calls, validates each CSV against the shipped schema, and refuses to
package one that does not match, because the organisers score the zip as-is. It then assembles the
team folder: \texttt{predictions.zip} flat at the top, the runnable app with its checkpoints, the
demo video, and \texttt{Optional\_Items/} holding this write-up and each subsystem's code and
model. \texttt{./scripts/deploy.sh} serves that same app on Cloud Run.

\section{Assumptions and limits}

\begin{itemize}
  \item Splits are by file, never by window, and every validation figure quoted above is
        out-of-fold. Nothing is tuned against the held-out test inputs.
  \item Rail: no fault file in training is slower than 9.70\,m/s while 133 of 234 \texttt{Normal}
        files are, so ``slow implies Normal'' is available for free. On the 138 files fast enough
        to remove that shortcut the score is \textbf{0.770}, above the headline rather than below
        it. We quote both figures rather than the flattering one.
  \item Rail: on this data, 36 columns of shuffled noise buy $+0.012$ macro F1, so any added
        feature block must beat its own shuffled control before it ships. Side I, 14 files, stays
        the weakest class at 0.515.
  \item ACV is a rule, not a fitted model, validated on five cases. Cheap to audit, but untested
        against a fleet carrying more than one fault.
  \item SHM damage values describe the damage contributed by individual recordings, not remaining
        asset life. The app's outputs are decision support for maintenance review, not a
        fitness-for-service decision.
\end{itemize}

\end{document}
